\documentclass[11pt,a4paper]{article}
\usepackage[utf8]{inputenc}
\usepackage[T1]{fontenc}
\usepackage{tgpagella}
\usepackage[spanish]{babel}
\usepackage{geometry}
\usepackage{booktabs}
\usepackage{xcolor}
\usepackage{titlesec}
\usepackage{fancyhdr}
\usepackage[parfill]{parskip}
\usepackage[expansion=false]{microtype}
\usepackage{hyperref}
\usepackage{setspace}
\usepackage{needspace}
\usepackage{etoolbox}
\usepackage{graphicx}

\widowpenalty=10000
\clubpenalty=10000
\displaywidowpenalty=10000
\raggedbottom
\sloppy
\hbadness=10000
\hfuzz=2pt

\geometry{top=3.0cm,bottom=3.0cm,left=3.5cm,right=3.5cm}
\setlength{\parskip}{0.65em}
\setlength{\parindent}{0em}
\setlength{\headheight}{14pt}

\definecolor{azulai}{RGB}{30,80,140}
\definecolor{doradoai}{RGB}{140,90,0}
\definecolor{grisai}{RGB}{70,70,70}

\titleformat{\section}{\Large\bfseries\color{azulai}}{}{0em}{}[{\color{azulai}\titlerule[0.5pt]}]
\titlespacing*{\section}{0pt}{1.8em}{0.8em}

\titleformat{\subsection}{\large\bfseries\color{doradoai}}{}{0em}{}
\titlespacing*{\subsection}{0pt}{1.4em}{0.5em}

\titleformat{\subsubsection}{\normalsize\bfseries\color{grisai}}{}{0em}{}
\titlespacing*{\subsubsection}{0pt}{1.0em}{0.3em}

% Evita títulos huérfanos al final de página
\pretocmd{\section}{\Needspace{12\baselineskip}}{}{}
\pretocmd{\subsection}{\Needspace{9\baselineskip}}{}{}
\pretocmd{\subsubsection}{\Needspace{6\baselineskip}}{}{}

\pagestyle{fancy}
\fancyhf{}
\rhead{\textcolor{grisai}{\small Virginia Riezu — Arquitectura Interna}}
\lhead{\textcolor{grisai}{\small Luna · Casa 4 · Casa 6}}
\cfoot{\textcolor{grisai}{\small\thepage}}
\renewcommand{\headrulewidth}{0.3pt}

\hypersetup{colorlinks=true,linkcolor=azulai,urlcolor=azulai}
\setstretch{1.45}
\tolerance=2500
\emergencystretch=8em

\begin{document}

% ── Portada ──────────────────────────────────────────────────────────────────
\begin{titlepage}
  \centering
  \vspace*{1.5cm}
  {\Huge\bfseries\color{azulai} Luna · Casa 4 · Casa 6}\\[0.5cm]
  {\large\color{grisai} Arquitectura Interna}\\[0.3cm]
  {\small\itshape\color{grisai} Un método para sostener cuerpo, energía y vida con coherencia}\\[2cm]
  {\huge\color{doradoai} Virginia Riezu}\\[1.5cm]
  {\Large 06/09/1975 \quad 01:30}\\[0.3cm]
  {\Large Pamplona, España}\\[0.3cm]
  {\normalsize Lat: 42.8157° \quad Lon: -1.6522° \quad UTC+02:00}\\[0.3cm]
  {\normalsize Ascendente: Géminis 29°04' \quad
    MC: Piscis 3°22'}\\[2cm]

  \begin{tabular}{lll}
    \textbf{Luna:} & Virgo & Casa 4 \\
    \textbf{Casa 4:} & Virgo en cúspide & Regente: Mercurio \\
    \textbf{Casa 6:} & Escorpio en cúspide & Regente: Plutón \\
  \end{tabular}\\[2cm]

  \vfill
  {\small Generado el 20/05/2026}
\end{titlepage}

\tableofcontents
\clearpage

% ── Datos de referencia ───────────────────────────────────────────────────────
\section{Datos de referencia}

\begin{center}
\begin{tabular}{llll}
  \toprule
  \textbf{Eje} & \textbf{Signo} & \textbf{Posición} & \textbf{Datos} \\
  \midrule
  Luna & Virgo & Casa 4 & 15°17' \\
  Casa 4 (cúspide) & Virgo & --- & Regente: Mercurio \\
  Casa 6 (cúspide) & Escorpio & --- & Regente: Plutón \\
  Ascendente & Géminis & --- & 29°04' \\
  Medio Cielo & Piscis & --- & 3°22' \\
  \bottomrule
\end{tabular}
\end{center}

\vspace{0.4cm}
\textbf{Planetas en Casa 4:} Sol, Luna\\
\textbf{Planetas en Casa 6:} Neptuno

\vspace{0.5cm}
\textbf{Aspectos de la Luna:}

\begin{center}
\begin{tabular}{lllll}
  \toprule
  & \textbf{Asp.} & \textbf{Planeta} & \textbf{Tipo} & \textbf{Orbe} \\
  \midrule
  Luna & □ & Marte & Cuadratura & 2.4° \\
  Luna & = & Sol & Conjunción & 2.5° \\
  Luna & □ & Neptuno & Cuadratura & 6.2° \\
  Luna & ☍ & Lilith & Oposición & 8.5° \\
  \bottomrule
\end{tabular}
\end{center}

\vspace{0.8cm}
\begin{center}
\includegraphics[width=0.80\textwidth]{C:/Users/Izas/claude/astrología/pruebas/Virginia_Riezu_Luna_C4_C6_rueda.png}
\end{center}


\vspace{0.5cm}
\Needspace{5\baselineskip}

% ── Interpretación ────────────────────────────────────────────────────────────
\section{Interpretación — Arquitectura Interna}

\Needspace{8\baselineskip}
\begin{center}
{\small\itshape
No se trata de explicarte. Se trata de mostrar cómo funciona tu regulación interna,\\
dónde puede haber desregulación y qué necesitas para sostenerte.
}
\end{center}

\vspace{0.8cm}
\Needspace{5\baselineskip}

% ── 1. Luna ───────────────────────────────────────────────────────────────────
\subsection{1. Tu Luna — Función emocional}

\subsubsection*{Luna en Virgo — Casa 4}

Tiendes a analizar lo que sientes antes de terminar de sentirlo.

Tu mente intenta entender rápidamente si la emoción tiene sentido, si es razonable o si debería estar ocurriendo. Muchas veces, lo que no consigues procesar emocionalmente termina desplazándose hacia el cuerpo.

La tensión suele aparecer entonces como cansancio, sobrecarga mental o sensación de no llegar a todo. También puede existir una tendencia constante a corregirte internamente sin darte cuenta.

La regulación mejora cuando existe orden simple, estructura cotidiana y atención concreta al cuerpo. En cambio, el exceso de exigencia, el sobreanálisis o la sensación de que nunca es suficiente tienden a aumentar la presión interna.

Tu mundo emocional está profundamente unido a la sensación de hogar y pertenencia. Los cambios familiares, las tensiones en casa o la pérdida de referencias afectivas te afectan mucho más de lo que suele verse desde fuera. Cuando algo te descoloca, tiendes a ir hacia dentro: buscar refugio, aislarte un poco o volver a lo conocido para recuperar sensación de seguridad. Necesitas sentir que existe un lugar donde puedes bajar la guardia de verdad. Cuando esa sensación de refugio falta, el cuerpo permanece mucho más tiempo en tensión. Te ayuda la intimidad, la estabilidad emocional y sentir que tienes una base a la que volver. En cambio, la tensión familiar, la pérdida de seguridad afectiva o sentir que no tienes refugio suele hacer que aparezca mucha desprotección interna.

Cuando algo emocionalmente te activa, el cuerpo entra rápidamente en tensión. Si no existe una forma clara de descargar esa energía, puede aparecer irritabilidad, impulsividad o necesidad de confrontación. Muchas veces la reacción responde más al nivel de activación acumulada que a lo que realmente está ocurriendo en ese momento. El aprendizaje no consiste en dejar de sentir intensidad, sino en aprender a darle una salida y un canal más consciente.

Lo que sientes y la dirección de tu vida suelen moverse muy unidas. Cuando emocionalmente hay estabilidad, aparece también más energía, claridad y sensación de propósito. Y cuando algo se rompe por dentro, la vitalidad suele notarlo enseguida. No separas fácilmente quién eres de cómo te sientes en cada momento. Eso puede darte mucha coherencia interna, pero también hacer que los períodos difíciles te atraviesen con más profundidad.

A veces cuesta distinguir claramente lo que sientes de lo que querrías sentir o imaginar sentir. Cuando emocionalmente necesitas claridad, puede aparecer todavía más confusión, mezcla o ambigüedad interna. La percepción emocional puede llegar mezclada con expectativas, idealización o imágenes internas difíciles de separar de la realidad concreta. Necesitas mucho más suelo emocional y referencia clara de lo que normalmente crees. Tu regulación mejora cuando existe descanso, límites y contacto real con lo que está ocurriendo aquí y ahora.

Puede haber oscilación entre buscar cuidado y rechazarlo cuando se siente demasiado invasivo o absorbente. Una parte de ti necesita pertenecer, sentirse acogido y poder descansar emocionalmente en alguien o en algo. Y al mismo tiempo existe otra parte que se rebela rápidamente ante cualquier forma de dependencia emocional o sensación de pérdida de libertad interna. Tu regulación mejora cuando puedes reconocer ambas necesidades sin expulsar ninguna de ellas.

\vspace{0.8cm}
\Needspace{5\baselineskip}

% ── 2. Casa 4 ─────────────────────────────────────────────────────────────────
\subsection{2. Tu Casa 4 — Raíz interna}

\subsubsection*{Virgo en la cúspide — Regente: Mercurio}

Cuando aparece sensación de desorden dentro de ti, muchas veces intentas recuperar estabilidad organizando algo fuera. Ordenar, limpiar, resolver tareas o estructurar el entorno puede ayudarte a bajar momentáneamente la tensión interna. La dificultad aparece cuando pasas demasiado tiempo intentando corregir lo externo sin terminar de entrar en lo que realmente te está ocurriendo por dentro. Tu sensación de seguridad aumenta cuando la vida cotidiana tiene cierto orden y funcionalidad. Cuando todo alrededor se vuelve caótico, el cuerpo entra rápidamente en sobrecarga. Te ayuda sentir estructura simple, orden y claridad en lo cotidiano. En cambio, el caos, la sobreexigencia o sentir que nunca llegas a todo suele generar mucha tensión interna.

El regente de la Casa 4 es Mercurio, en Libra y en Casa 5. La base interna tiende a organizarse desde la creatividad y la expresión, con la cualidad mental y relacional de Libra.

Necesitas sentir que puedes ser plenamente quien eres dentro de tu espacio más íntimo. Cuando existe demasiada distancia entre lo que muestras hacia fuera y lo que realmente sucede dentro de ti, algo empieza a desgastarse lentamente. No suele ayudarte sostener una imagen que no tiene raíces reales en tu vida cotidiana, porque tu mundo interno necesita coherencia y verdad emocional. El hogar no funciona solo como un lugar físico. Es el espacio donde deberías poder bajar la exigencia y existir sin tener que sostener constantemente una versión de ti que no se siente auténtica. Cuando no existe ese lugar de autenticidad y descanso interno, la energía empieza a agotarse silenciosamente.

Tu mundo emocional está profundamente unido a tus raíces y al entorno donde vives. Cuando hay tensión en casa o inestabilidad emocional alrededor, recuperar calma se vuelve mucho más difícil. Necesitas intimidad, refugio y cierta sensación de seguridad emocional para descansar de verdad. Los cambios en el hogar no te afectan solo de forma práctica. También se mueven profundamente por dentro y el cuerpo suele notarlo muy rápido cuando algo altera tu espacio emocional. Cuando el entorno íntimo pierde estabilidad, todo dentro necesita mucho más tiempo para volver a sentirse seguro.

\vspace{0.8cm}
\Needspace{5\baselineskip}

% ── 3. Casa 6 ─────────────────────────────────────────────────────────────────
\subsection{3. Tu Casa 6 — Regulación cotidiana}

\subsubsection*{Escorpio en la cúspide — Regente: Plutón}

Tu cuerpo funciona mucho por acumulación y descarga. Cuando acumulas demasiada intensidad sin darle salida, la tensión empieza a quedarse dentro. Las prácticas demasiado suaves muchas veces no son suficientes para ti. Necesitas descargar de verdad: sudar, atravesar intensidad o entrar en períodos de silencio profundo. Cuando no existe esa descarga, el cuerpo empieza a endurecerse y la energía queda retenida durante demasiado tiempo. Entonces pueden aparecer tensión constante, dificultad para dormir o sensación de pesadez acumulada. Te ayudan la descarga profunda, la intensidad bien canalizada y los espacios de silencio real. En cambio, acumular emociones, contener demasiado tiempo lo que sientes o no tener salida para ello suele generar muchísima presión interna.

El regente de la Casa 6 es Plutón, en Libra y en Casa 5. La regulación cotidiana tiende a organizarse desde la creatividad y la expresión, con la cualidad mental y relacional de Libra.

Muchas veces te cuesta notar claramente dónde termina el esfuerzo y dónde empieza el agotamiento. Puedes absorber la carga emocional del entorno sin darte cuenta enseguida y seguir funcionando mientras el cuerpo se va saturando lentamente. Necesitas espacios reales de descanso, silencio y cierre emocional para vaciar todo lo que vas acumulando. Cuando no existe ese vaciado, la fatiga empieza a crecer aunque no haya una causa física evidente. Tu cuerpo necesita mucha más recuperación emocional de la que suele parecer desde fuera.

\vspace{0.8cm}
\Needspace{5\baselineskip}

% ── 4. Integración ────────────────────────────────────────────────────────────
\subsection{4. Integración — Cómo funcionan las tres capas juntas}

Tu Luna en Virgo muestra cómo reaccionas emocionalmente cuando algo te afecta.

Es el primer movimiento: lo que aparece antes de que puedas ordenarlo, explicarlo o decidir qué hacer con ello.

Después, esa emoción busca un lugar donde apoyarse. Ahí entra tu Casa 4, con Virgo en la cúspide: la parte de ti que necesita sentir base, refugio y seguridad interna.

Si esa base está disponible, lo que sientes puede moverse con más facilidad.
Si no lo está, la emoción se queda más tiempo en estado de alerta.

Y lo que no encuentra apoyo interno termina llegando al cuerpo. Ahí entra tu Casa 6, con Escorpio en la cúspide: tus hábitos, tu energía diaria, tu descanso y la forma en que el cuerpo sostiene lo que vives.

Por eso, cuando algo se desordena emocionalmente, no se queda solo en lo emocional. Puede acabar apareciendo como cansancio, pérdida de ritmo, dificultad para descansar o sensación de no poder sostener el día con la misma energía.

Tu Luna y tu Casa 4 comparten el elemento Tierra.

Eso significa que lo que sientes y lo que necesitas para sentir base interna tienen un lenguaje parecido.

Cuando emocionalmente algo se activa, una parte profunda de ti reconoce bastante rápido qué necesita para sentirse más segura.

La Casa 6, en cambio, está en Agua. Tu cuerpo y tus hábitos funcionan con otra lógica.

Puede que entiendas emocionalmente lo que necesitas, pero aun así te cueste llevarlo al ritmo cotidiano, al descanso, al cuerpo o a la rutina.

Además, el Sol y la Luna en la Casa 4 influye directamente en tu sensación de base interna.

Cuando buscas seguridad, refugio o estabilidad emocional profunda, esa energía ya está presente ahí.

Por eso tu forma de sentir hogar, intimidad y raíz no depende solo del signo de la Casa 4. También está marcada por lo que esos planetas activan dentro de ti.

Además, Neptuno en la Casa 6 influye en cómo tu cuerpo sostiene la vida cotidiana.

No se trata solo de hábitos o rutina. Se trata de cómo organizas tu energía, cómo respondes al cansancio y qué ocurre cuando el cuerpo empieza a pedir ajuste.

Esos planetas muestran qué tipo de fuerza, sensibilidad o tensión entra en tu regulación diaria.

Los aspectos de tensión con Marte, Neptuno y Lilith hacen que el mundo emocional tenga menos margen en ciertos momentos.

Cuando algo te activa, puede haber menos tiempo entre sentir y reaccionar. La emoción puede llegar más intensa, más mezclada o más difícil de ordenar.

Esto no significa que sí o sí vayas a desregularte. Significa que necesitas reconocer antes las señales iniciales, porque cuando la activación ya ha subido mucho, cuesta más volver al centro.

Al mismo tiempo, los aspectos de apoyo con el Sol ofrecen recursos reales.

Cuando consigues encontrar un primer punto de estabilidad, esos aspectos ayudan a sostenerlo.

No eliminan la tensión, pero pueden facilitar que vuelvas a organizarte sin quedarte completamente dentro de lo que se activó.

Lo más difícil no suele ser que una sola parte se altere.

Lo más difícil aparece cuando las tres capas se mueven al mismo tiempo:
la Luna en Virgo, con su forma concreta de sentir;
la Casa 4 en Virgo, con su necesidad de base y seguridad interna;
y la Casa 6 en Escorpio, con el cuerpo intentando sostenerlo todo en la vida diaria.

Cuando esto ocurre, puedes notar que no sabes por dónde empezar: emocionalmente hay activación, internamente falta suelo y el cuerpo empieza a mostrar cansancio, tensión o pérdida de ritmo.

El trabajo no consiste en resolverlo todo a la vez.

Consiste en encontrar la primera capa disponible: una emoción que pueda expresarse, un entorno que pueda darte más seguridad, o una acción corporal sencilla que te ayude a volver poco a poco a la estabilidad.

\vspace{0.8cm}
\Needspace{5\baselineskip}

% ── 5. Sostén y desregulación ─────────────────────────────────────────────────
\subsection{5. Sostén y desregulación}

Tu Luna en Virgo se desregula cuando lo que sientes no encuentra salida.

Puede ser porque no puedes expresarlo, porque no hay espacio para moverte, porque el entorno no lo recibe o porque intentas sostenerlo demasiado tiempo dentro.

No siempre hace falta que ocurra algo muy intenso.

A veces basta con acumular durante demasiado tiempo una emoción que no ha tenido lugar, palabra, movimiento o descanso.

Los aspectos de tensión con Marte, Neptuno y Lilith hacen que ese margen sea más estrecho.

Cuando algo te activa, puede haber menos tiempo entre sentir y reaccionar.

La emoción puede llegar más fuerte, más mezclada o más difícil de ordenar. Por eso es importante reconocer las primeras señales antes de llegar al punto de desborde.

Tu Casa 4 con Virgo pierde estabilidad cuando dejas de sentir base interna.

No siempre tiene que haber un conflicto claro.

A veces basta con que cambie el clima emocional del entorno, con que una relación cercana se tense, con que no puedas descansar de verdad o con que el espacio íntimo deje de sentirse seguro.

Cuando eso ocurre, no solo cambia el ánimo. Cambia la sensación de tener suelo por dentro.

Tu Casa 6 con Escorpio muestra la desregulación en el cuerpo y en la vida cotidiana.

Puede aparecer como cansancio, pérdida de ritmo, tensión física, sueño menos reparador o dificultad para sostener hábitos que normalmente sí puedes sostener.

Muchas veces el cuerpo avisa antes de que puedas explicar qué está pasando.

Cuando el cuerpo empieza a perder continuidad, conviene mirar también qué emoción no ha encontrado salida y qué parte de tu base interna ha dejado de sentirse segura.

Los momentos más difíciles aparecen cuando las tres capas se mueven a la vez.

La Luna en Virgo se activa.
La Casa 4 en Virgo deja de sentirse como base.
Y la Casa 6 en Escorpio empieza a mostrarlo en el cuerpo, el descanso o la rutina.

Ahí puedes sentir que no sabes por dónde empezar.

No se trata de resolverlo todo a la vez.
Se trata de encontrar la primera capa disponible: expresar algo, descansar, ordenar una pequeña rutina, moverte, pedir espacio o recuperar una mínima sensación de seguridad.

Los aspectos de apoyo con el Sol funcionan como recursos cuando consigues recuperar un primer punto de estabilidad.

No eliminan la desregulación, pero pueden ayudarte a reorganizarte antes.

Cuando aparece una mínima base, esos aspectos facilitan que no te quedes dentro de lo que se activó.

\vspace{0.8cm}
\Needspace{5\baselineskip}

% ── 6. Orientación final ──────────────────────────────────────────────────────
\subsection{6. Orientación final}

\subsubsection*{Qué observar}

Empieza observando qué situaciones activan tu Luna en Virgo.

No para corregirte.
No para controlar lo que sientes.

Solo para reconocer antes cómo empieza la activación.

¿Qué ocurre justo antes de sentirte que la situación te sobrepasa?
¿Qué cambia en el entorno, en los vínculos o en el ritmo diario?

Muchas veces la desregulación no empieza en el momento del desborde. Empieza bastante antes: cuando tu Casa 4 en Virgo deja de sentirse segura o cuando tu Casa 6 en Escorpio pierde continuidad y el cuerpo empieza a acumular tensión.

\subsubsection*{Qué puede desregularte}

En tu caso, la desregulación suele seguir una secuencia.

Primero, tu Luna en Virgo busca algo que necesita para estabilizarse y no lo encuentra.

Después, tu Casa 4 deja de funcionar como base interna. Puede ser por tensión emocional, por cambios en el entorno, por exceso de carga o simplemente porque llevas demasiado tiempo sosteniendo más de lo que puedes absorber.

Y finalmente, la Casa 6 en Escorpio empieza a mostrarlo en el cuerpo:
menos energía,
menos descanso,
menos capacidad de sostener hábitos o continuidad.

Cuando las tres capas coinciden, puedes sentir que todo cuesta más de lo habitual aunque desde fuera no parezca estar ocurriendo nada tan grave.

\subsubsection*{Qué te estructura y sostiene}

Lo que te ayuda a regularte no tiene que ser enorme.

De hecho, muchas veces lo más importante son cosas pequeñas que se sostienen en el tiempo.

Una mínima rutina corporal.
Un espacio de descanso real.
Un entorno que se sienta más seguro.
Una forma de expresar lo que ocurre antes de acumularlo demasiado.

Cuando tu Casa 4 siente suelo, tu Luna necesita menos esfuerzo para estabilizarse.

Y cuando tu Casa 6 mantiene cierta continuidad —aunque sea pequeña— el cuerpo puede absorber mejor lo que emocionalmente estás atravesando.

No se trata de no sentir.
Se trata de que lo que sientes no tenga que terminar descargándose siempre en agotamiento, tensión o pérdida de equilibrio.

\subsubsection*{Práctica}

Tu regulación necesita vaciado emocional.

Dedica un momento del día a cerrar lo absorbido:
silencio,
respiración,
una ducha consciente,
música tranquila,
o simplemente estar sin estímulos.

No es una pausa superficial.
Es una forma de ayudarte a soltar lo que has ido acumulando.

Cuando no existe ese cierre, el cuerpo sigue cargando emociones incluso mientras intenta descansar.

\Needspace{8\baselineskip}
\vspace{0.6cm}

\begin{center}
{\small\itshape\color{grisai}
La astrología se usa aquí como lenguaje simbólico de observación, no como definición de la persona.\\
Lo que tienes delante es un mapa de posibilidades, no un destino fijo.
}
\end{center}

\end{document}
